\subsection{导子的扩张：域的情形}

设 K、L 和 M 为域，使得 \(K \subset L \subset M\) 。由命题 21（III，第 572 页），我们得到向量 M-空间的一个正合序列

\[
(E _ {K, L, M})
\]

\[
\Omega_ {\mathrm{K}} (\mathrm{L}) \otimes_ {\mathrm{L}} \mathbf {M} \xrightarrow {\alpha} \Omega_ {\mathrm{K}} (\mathbf {M}) \xrightarrow {\beta} \Omega_ {\mathrm{L}} (\mathbf {M}) \to 0;
\]

M-线性映射 \(\alpha\) 和 \(\beta\) 由以下关系刻画

\begin{enumerate}
\def\labelenumi{(\arabic{enumi})}
\setcounter{enumi}{2}
\tightlist
\item
\end{enumerate}

\[
\begin{array}{r l} \alpha (d _ {\mathrm{L/K}} x \otimes 1) = d _ {\mathrm{M/K}} x & \text { for } \quad x \in \mathrm{L} \\ \beta (d _ {\mathrm{M/K}} y) = d _ {\mathrm{M/L}} y & \text { for } \quad y \in \mathrm{M}. \end{array}\tag{4}
\]

设 V 是一个向量 M-空间，记 \(D_{K}(M,V)\) 为 M 的取值于 V 的 K-导子构成的向量空间，并类似地引入 \(D_{K}(L,V)\) 和 \(D_{L}(M,V)\) 。由 III，第 571 页，图表 (42) 和 III，第 572 页，图表 (44)，我们有一个向量空间同态的交换图

\[
\begin{array}{l}0 \rightarrow \operatorname{Hom} _ {\mathbf {M}} (\Omega_ {\mathrm{L}} (\mathbf {M}), \mathrm{V}) \xrightarrow {\operatorname{Hom} (\beta , 1)} \operatorname{Hom} _ {\mathbf {M}} (\Omega_ {\mathrm{K}} (\mathbf {M}), \mathrm{V}) \xrightarrow {\operatorname{Hom} (\alpha , 1)} \operatorname{Hom} _ {\mathbf {M}} (\Omega_ {\mathrm{K}} (\mathrm{L}) \otimes_ {\mathrm{L}} \mathrm{M}, \mathrm{V})\\\downarrow\\0 \rightarrow \operatorname{Der} _ {\mathrm{L}} (\mathrm{M}, \mathrm{V}) \xrightarrow {i _ {\mathrm{V}}} \operatorname{Der} _ {\mathrm{K}} (\mathrm{M}, \mathrm{V}) \xrightarrow {r _ {\mathrm{V}}} \operatorname{Der} _ {\mathrm{K}} (\mathrm{L}, \mathrm{V}),\end{array}
\]

其中竖直箭头是同构， \(i_{v}\) 是典范嵌入，而 \(r_{v}\) 是限制映射。

因此，由 II，第 299 页，定理 5 和 II，第 301 页，命题 10，我们得到以下命题：

命题 2. --- 以下条件等价：

\begin{enumerate}
\def\labelenumi{\alph{enumi})}
\item
  映射 \(\alpha\) 是单射的。
\item
  L 到 M 的每个 K-导子都延拓为 M 到 M 的一个 K-导子。
\item
  L 到向量 M-空间 V 的每个 K-导子都延拓为 M 到 V 的一个 K-导子。
\end{enumerate}

命题 3. --- 设 K 是一个域，L 是 K 的一个纯扩张，且 \((x_{i})_{i\in I}\) 是 L 的一个纯基（V，第 106 页，定义 2）。

\begin{enumerate}
\def\labelenumi{\alph{enumi})}
\item
  设 \(\mathbf{V}\) 是 \(\mathbf{L}\) 上的向量空间， \(\Delta\) 是 \(\mathbf{K}\) 到 \(\mathbf{V}\) 中的一个导子，且 \((v_i)_{i\in I}\) 是 \(\mathbf{V}\) 的元素的一个族。存在唯一导子 \(\mathbf{D}\) （\(\mathbf{L}\) 到 \(\mathbf{V}\) 中的），它延拓 \(\Delta\) 且满足 \(\mathbf{D}(x_i) = v_i\) 对所有 \(i\in I\) 。
\item
  L 的 K-微分的向量 L-空间 \(\Omega_{\mathbf{K}}(\mathbf{L})\) 以族 \((dx_i)_{i\in I}\) 为基。
\end{enumerate}

断言 b) 已在 IV，第 23 页得到证明，断言 a) 直接由此得出。

推论. --- 设 P 是 K 的一个子域。典范映射 \(\alpha\) （\(\Omega_{\mathrm{P}}(\mathrm{K}) \otimes_{\mathrm{K}} \mathrm{L}\) 到 \(\Omega_{\mathrm{P}}(\mathrm{L})\) 中的）是单射，且族 \((d_{\mathrm{L}/\mathrm{P}}x_{i})_{i \in \mathrm{I}}\) 是 \(\Omega_{\mathrm{P}}(\mathrm{L})\) 的一个子空间（在 L 上的）的一组基，该子空间是 \(\alpha: \Omega_{\mathrm{P}}(\mathrm{K}) \otimes_{\mathrm{K}} \mathrm{L} \to \Omega_{\mathrm{P}}(\mathrm{L})\) 的像的补空间。

命题 3 的 a) 表明，K 到 L 上向量空间 V 的每一个 P-导子都延拓为 L 到 V 的 P-导子； \(\alpha\) 的单射性于是由命题 2 得出。推论的第二个断言由序列 \((E_{P,K,L})\) 的正合性以及命题 3, b)（考虑到公式 (4)）得出。

命题 4. --- 设 K 是一个域，L 是 K 的一个可分代数扩张，V 是 L 上的一个向量空间。

\begin{enumerate}
\def\labelenumi{\alph{enumi})}
\item
  L 到 V 的每一个 K-导子都为零。
\item
  如果 \(\Delta\) 是 \(K\) 到 \(V\) 中的一个导子，则存在唯一的导子 \(D\) （\(L\) 到 \(V\) 中的），它延拓 \(\Delta\) 。
\end{enumerate}

设 D 为 L 到 V 的一个 K-导子。若 E 是 L 的一个在 K 上有限次的子扩张，则 K-代数 E 是平展的，因此 \(\Omega_{\mathrm{K}}(\mathrm{E}) = 0\) （V，第 33 页，定理 3），由此 \(D \mid E = 0\) （根据 \(\Omega_{\mathrm{K}}(\mathrm{E})\) 的泛性质，III，第 569 页）。由于 L 是 K 上有限次子扩张的一个族的并，我们有 D = 0，从而得出 a)。

设 \(\Delta\) 为 K 到 V 的一个导子。若 D' 和 D'' 是 \(\Delta\) 到 L 到 V 的导子的两个延拓，则差 D' - D'' 是 L 到 V 的一个 K-导子；因此由 \(a\) ) 可知它为零，所以 D' = D''。

仍需证明 \(\Delta\) 的一个延拓的存在性。Zorn 引理（E, III, 第 20 页）蕴含存在一个极大延拓 \(D_0\) （\(\Delta\) 的延拓），它是定义在子域 \(L_0\) （\(L\) 的、包含 \(K\) 的）上的导子。

设 \(x\) 属于 \(L\) ， \(g\) 为 \(x\) 在 \(L_0\) 上的极小多项式。由于 \(L\) 在 \(K\) 上是代数的且可分的， \(x\) 在 \(L_0\) 上是代数的且可分的（V，第 39 页，命题 6 及第 39 页，推论 2）；因此 \(x\) 是 \(g\) 的单根（V，第 39 页，命题 5），由此得 \(g'(x) \neq 0\) （IV，第 17 页，命题 7）。如果我们定义 \(g^{D_0}(x)\) 如 \(V\) ，第 127 页中那样，则存在元素 \(u\) （\(V\) 的）使得 \(g^{D_0}(x) + g'(x) \cdot u = 0\) ；根据命题 1（V，第 127 页），存在一个导子 \(D\) （\(L_0(x)\) 到 \(V\) 中的），它延拓 \(D_0\) 且满足 \(D(x) = u\) 。鉴于 \((L_0, D_0)\) 的极大性，我们于是有 \(L_0(x) = L_0\) ，从而 \(x \in L_0\) 。由于 \(x\) 是任意的，我们得出结论 \(L_0 = L\) 。

推论 1. --- 我们有 \(\Omega_{\mathbf{K}}(\mathbf{L}) = 0\) ，如果 \(\mathbf{L}\) 在 \(\mathbf{K}\) 上是代数的且可分的。

该推论由命题 4 的 a) 得出，因为向量 L-空间 \(\Omega_{\mathbf{K}}(\mathbf{L})\) 由典范 K-导子 \(d_{\mathrm{L / K}}: \mathrm{L} \to \Omega_{\mathbf{K}}(\mathrm{L})\) 的像生成。

推论 2. --- 若 L 是域 K 的一个代数可分扩张，则典范映射 \(\alpha: \Omega_{\mathrm{P}}(\mathrm{K}) \otimes_{\mathrm{K}} \mathrm{L} \to \Omega_{\mathrm{P}}(\mathrm{L})\) 对 K 的每一个子域 P 都是同构。

映射 \(\alpha\) 由命题 2（V，第 128 页）和命题 4, \(b\) ) 是单射的；由于 \(\Omega_{\mathrm{K}}(\mathrm{L})\) 化为 0（推论 1），序列 \((\mathrm{E}_{\mathrm{P},\mathrm{K},\mathrm{L}})\) 的正合性蕴含 \(\alpha\) 是满射。

推论 3. --- 设 E 是域 K 的一个扩域，D 是 E 到 E 的一个导子，且将 K 映入 K。若 L 是 E 的一个子扩张，在 K 上是代数的且可分的，则 D(L) ⊂ L。

设 \(\Delta\) 为 K 到 L 的导子，它在 K 上与 D 一致。由命题 4（V，第 129 页），存在 L 到 L 的一个导子 \(D'\) ，它延拓 \(\Delta\) 。现在我们可以把 \(D'\) 以及 D 在 L 上的限制 \(D''\) 都看作 L 到 E 的导子；由于它们在 K 上一致，由命题 4 我们有 \(D' = D''\) ，由此得到

\[
\mathrm{D} (\mathrm{L}) = \mathrm{D} ^ {\prime \prime} (\mathrm{L}) = \mathrm{D} ^ {\prime} (\mathrm{L}) \subset \mathrm{L}.
\]

注记. --- 1) 稍后（V，第 131 页，推论 3 及第 135 页，推论 2）我们将证明命题 4 的推论 1 的一个逆命题。

\begin{enumerate}
\def\labelenumi{\arabic{enumi})}
\setcounter{enumi}{1}
\tightlist
\item
  素域的每个代数扩张都是可分的（V，第 37 页，推论）。由于素域的每个导子都是零，素域的代数扩张的每个导子也都是零（命题 4）。
\end{enumerate}

